\chapter{Claim-to-proof and evidence register}
\label{sec:towers-claim-map}
The register covers the mathematical statements, rounded evaluations and
nonmathematical assertions in the article and its audit appendices. A model
assumption is an input; an observed execution or a historical attribution is
an evidence claim. Neither can be turned into a theorem about actual tests.
The exact expressions, rather than the printed decimal digits, are what the
proofs establish. The bibliography identifies the well-known results used;
their relevant deductions are proved in the preceding appendices.

\small
\begin{longtable}{>{\raggedright\arraybackslash}p{.28\linewidth}>{\raggedright\arraybackslash}p{.65\linewidth}}
\toprule
Claim or location & Proof, exact specification, or evidence boundary\\
\midrule\endhead
Abstract and source audit; Table~\ref{tab:original-audit}
& Rounded values evaluate \eqref{eq:joint}, \eqref{eq:union} and
\eqref{eq:conditionals}; Propositions~\ref{tw:gaussian}--\ref{tw:events}
prove those formulas. Proposition~\ref{tw:original-bound} rigorously rejects
the principal first-letter probabilities. The source columns, retraction,
dates and Cattell recollection are attributed records
\citep{towersOriginal,towersFollowup}; Cattell's original page was not
inspected, and no theorem proves an attribution.\\
Joint model \eqref{eq:model}; cutoffs \eqref{eq:cutoffs}
& Gaussianity and \(\rho=0.7\) are assumptions. Proposition~\ref{tw:gaussian}
proves the model's generator, density, unique quantile cutoffs and null
boundary events. Proposition~\ref{pf:normal-marginals} gives a counterexample
to inference of joint Gaussianity from normal margins and correlation.\\
Joint, union, conditionals and phi; \eqref{eq:joint}--\eqref{eq:phi}
& Propositions~\ref{tw:gaussian}, \ref{tw:events}; all four cell masses,
positive denominators, complements and indicator moments are derived.\\
Equal-cutoff results; \eqref{eq:equal-results}--\eqref{eq:union-denominator};
Table~\ref{tab:equal-cutoffs}
& Substitute \(p_A=p_B\) in Proposition~\ref{tw:events}; specified-only,
exactly-one and union-conditioned probabilities have distinct formulas.
IQ cutoffs use \(100+16\Phi^{-1}(1-p)\). Table entries are approximate
quadrature/quantile evaluations, not exact decimal equalities.\\
Mixed results \eqref{eq:mixed-joint}--\eqref{eq:mixed-remainder};
correlation sensitivity
& Propositions~\ref{tw:events}, \ref{tw:plackett}, \ref{tw:nesting} prove
conditional asymmetry, its factor twenty, correlation monotonicity and the
same-score versus different-test distinction. Displayed percentages are
rounded evaluations at the stated three correlations.\\
Rarity \eqref{eq:rarity}; calibrated rule \eqref{eq:calibration}
& Propositions~\ref{tw:rarity}, \ref{tw:calibration} prove quantile mapping,
explicit nonidentifiability counterexamples, and unique calibration. The
numerical root approximates the exact unique solution. No latent-ability
measurement relation is assumed or inferred.\\
Population, applicants, membership, retesting; \eqref{eq:bounds}
& Proposition~\ref{tw:attempts} proves the union bounds and independent
attempt formula, and supplies a selection counterexample. Actual testing
and membership claims require external evidence. Psychometric guidance,
range restriction and practice effects are attributed to
\citet{standards2014,sackett2000,hausknecht2007}; recommendations are stated
choices, not mathematical theorems.\\
BASIC ratio \eqref{eq:grid}; Tables~\ref{tab:grid}, \ref{tab:integer-cutoffs}
& Proposition~\ref{tw:grid} proves finite normalization, cancellation,
node count, midpoint boundary, quadrature bound and matching marginal
denominators. Published/reproduced digits and relative differences are
recorded approximate evaluations; the precise historical cause of their
remaining difference is unknown. Relative error means
\(\lvert x_{\mathrm{reproduced}}-x_{\mathrm{published}}\rvert/\lvert x_{\mathrm{published}}\rvert\).
The size comparison of numerical errors is numerical evidence.\\
Independent integral \eqref{eq:rho-integral}; numerical checks
& Proposition~\ref{tw:plackett} derives the density derivative, justifies
interchanging derivative and integral, and integrates from independence.
Symmetry and probability bounds follow in Propositions~\ref{tw:gaussian},
\ref{tw:events}. Agreement of algorithms is a numerical check, not a
certified integration-error proof.\\
Simulation cells, Wilson intervals and rare events in verification discussion
& Proposition~\ref{tw:simulation} proves exact bookkeeping and random
conditional denominators. Results~\ref{sm:binomial}, \ref{sm:wilson},
\ref{sm:clopper-pearson}, \ref{sm:rare} prove the ideal sampling laws and
interval construction; Wilson coverage is approximate, exact binomial
coverage conservative. No finite generator independence theorem is claimed.\\
Quantization and \(3.37\times10^{-11}\) bound
& Proposition~\ref{tw:quantization} derives the ratio perturbation bound
and its stated substitution. It does not bound integration or library error.\\
Python/Dask/Deal verification, numerical evidence and source identity
& \Cref{sec:verification} describes the maintained Python/Dask
reference calculations and Deal contracts. Numerical samples,
reproducible configuration, static lint and runtime/property
checks have specific scopes; they do not prove all executions,
external numerical libraries, random-generator independence
or population-model validity. Rocq checks separate mathematical
propositions with explicit hypotheses.\\
AI disclosure, provenance and limits
& \Cref{sec:ai-disclosure} records the project workflow and its
documented limits; it is not a probabilistic deduction or independent human
review. No theorem supplies missing empirical validation.\\
\bottomrule
\end{longtable}
\normalsize
